
This study set out to determine whether HuggingFace Hub dataset cards and OpenML run counts can serve as data sources for a study linking metadata-observable dataset misuse signals to ML reproducibility failure outcomes in Raff's 255-paper corpus. The infrastructure feasibility check (H-E1) reveals that they cannot: HF Hub achieves 30\% dataset card coverage (15/50 datasets found; all NLP benchmarks; large-scale DL vision, speech, and graph datasets absent), and OpenML achieves 22\% pre-publication temporal coverage (11/50; predominantly classical tabular datasets plus IMDB, 20 Newsgroups, and MNIST). Both are insufficient for valid regression analysis.

The coverage gaps are not random. They correspond precisely to each platform's founding community and architectural scope. HF Hub's ecosystem is NLP-centric; OpenML's architecture is tabular-centric. Any reproducibility auditing study of pre-2018 ML papers built on HF Hub and OpenML will be structurally biased toward NLP and classical ML papers, silently omitting the large-scale DL vision papers that dominate Raff's deep learning-era corpus.

This study's main contributions are: (1) the first empirical quantification of HF Hub and OpenML coverage for the Raff 2019 corpus, with exact coverage rates and domain breakdowns; (2) characterization of the domain-temporal selection bias in both APIs; (3) identification of Papers With Code as the most appropriate primary data source for a redesigned study; and (4) a validated 18-test data acquisition pipeline reusable for the redesigned study.

The original research question --- do metadata-observable dataset misuse signals predict ML reproducibility failure? --- remains open. The infrastructure gap reported here defines the methodological prerequisite for testing it. Any researcher planning to use HF Hub and OpenML to study pre-2018 ML reproducibility should explicitly measure API coverage for their target corpus before proceeding; the structural domain biases in these platforms are not visible without measurement.

